% =====================================================================
\section{Conclusion}
\label{sec:conclusion}
% =====================================================================

We asked whether the measured facts a deterministic supervisor writes into a
retry prompt are what repairs a defective stage in an LLM pipeline. The
pre-registered answer is no, for injected defects. A rejection that said
nothing repaired them as well as one that named the offending item and its
measured value: the difference was +0.7 percentage points, with an interval
that excludes anything larger than about 4. The reason is that the model,
shown its defective output, mostly copies it back and corrects the one
inconsistency it can see. The information the supervisor measures is in the
model's context too. A second model that repaired far less often gained an
exploratory 6.8 points from the measured facts, apparently because they told it
where the defect was; the null describes a capable model, not every model.

What the supervisor \emph{does} matters more than what it \emph{says}. Showing
the model its rejected output left repair unchanged and cut collateral damage
by an order of magnitude. The precision of a check decides whether a retry
fixes anything: on the defects the supervisor really caught in production, most
of them its own false positives, the production retry was accepted 28\% of the
time. Stage-local checks certify consistency with upstream, so a defect that
gets past its own stage is invisible to every later stage and is sometimes
rewarded by them. And a statistic the supervisor can compute is better
overwritten than retried. A retry also fixes it, but only an overwrite does so
exactly and at no cost.

For engineers, the lesson is to spend effort on where a check looks and how
precisely it fires, not on how its complaint is worded beyond naming what
failed. For researchers who
evaluate LLM systems by fault injection, the lesson is one that mutation
testing learned for detection: injected faults stand in for real ones only once
a coupling study shows that they do. Until then, a repair rate measured on
injected defects is a property of the operators. It is not yet evidence about
the defects the system actually produces.
